% rank-deficiency-dismissal.tex
%
% Lemma: KKT systems with a non-trivial null-space direction in the
% beta-block can be skipped.
% Intended location: between lem:well-defined and thm:algorithm-correct
% in general-case-algorithm-proof.tex.
%
% Sources: conjecture from Jörn (2026-02-19), proof by session.
% Jörn: math approved (8818fb8) — from \begin{lemma} to end of \begin{remark}

\begin{lemma}[Dismissal of redundant systems]%
\label{lem:rank-deficiency-dismissal}
Let $(S, \sigma)$ be a pair for which the linear
system~\eqref{eq:linear-system} has a solution
$(\beta_0, \lambda_0, \nu_0)$ with $\beta_0 \geq 0$,
and suppose the null space of the system matrix contains an element
$(\delta\beta, \delta\lambda, \delta\nu)$ with
$\delta\beta \neq 0$.

Then there exist $k \in S$ and a solution
$(\beta^*, \lambda^*, \nu^*)$ with
\begin{enumerate}
\item $\beta^*_k = 0$,
\item $\beta^*_j \geq 0$ for all $j \in S$,
\item $\beta^{*\top} H\, \beta^* = \beta_0^\top H\, \beta_0$.
\end{enumerate}
The restriction $\beta^*|_{S'}$ with $S' = S \setminus \{k\}$
lies in the constraint set $M(K)$ restricted to~$S'$,
so the algorithm for~$(S', \sigma|_{S'})$ maximises
$\beta^\top H'\, \beta$ over at least the value~$\nu_0$.
Consequently, the algorithm may discard any pair $(S, \sigma)$
whose system has such a null-space direction
without missing the capacity-achieving pair.
\end{lemma}

\begin{proof}
If $\beta_0$ already has $\beta_{0,k} = 0$ for some $k \in S$,
take $\beta^* = \beta_0$: conclusions (1)--(3) hold immediately
(with~(3) trivial), and the argument of Step~5 below gives the
restriction to~$S'$. So assume $\beta_{0,i} > 0$ for all~$i$.

Since $(\delta\beta, \delta\lambda, \delta\nu)$ is in the null space:
\[
  \text{(I$'$)}\ N^\top \delta\beta = 0, \qquad
  \text{(II$'$)}\ H\,\delta\beta
    = N\,\delta\lambda + \eta\,\delta\nu, \qquad
  \text{(III$'$)}\ \eta^\top \delta\beta = 0.
\]

\emph{Step~1: $\delta\beta$ has mixed signs.}
Expanding~(III$'$):
$\sum_{i \in S} \eta_i\, \delta\beta_i = 0$
where $\eta_i = h_{\sigma(i)} > 0$ are the heights.
Since $\delta\beta \neq 0$, at least one component is
nonzero, and the vanishing weighted sum with positive weights
forces $\delta\beta$ to have at least one strictly positive
and at least one strictly negative component.

\emph{Step~2: Walk along the null direction.}
Define $\beta(\alpha) = \beta_0 + \alpha\,\delta\beta$.
For every~$\alpha$, the triple
$(\beta(\alpha),\, \lambda_0 + \alpha\,\delta\lambda,\,
  \nu_0 + \alpha\,\delta\nu)$
satisfies~\eqref{eq:linear-system} (linearity of the system).
By Lemma~\ref{lem:well-defined}
(all solutions yield the same value of
$\beta^\top H\, \beta$),
$\beta(\alpha)^\top H\, \beta(\alpha) = \nu_0 + \alpha\,\delta\nu$.
On the other hand, expanding the quadratic form and using
$\delta\beta^\top H\, \beta_0 = 0$ and
$\delta\beta^\top H\, \delta\beta = 0$
(both from the proof of Lemma~\ref{lem:well-defined}):
\[
  \beta(\alpha)^\top H\, \beta(\alpha)
  = \beta_0^\top H\, \beta_0
    + 2\alpha\,\delta\beta^\top H\,\beta_0
    + \alpha^2\,\delta\beta^\top H\,\delta\beta
  = \nu_0.
\]
Comparing the two expressions: $\delta\nu = 0$, so the
multiplier~$\nu$ is the same for every solution.

\emph{Step~3: Identify the boundary point.}
Since $\beta_{0,i} > 0$ for all~$i$,
$\beta(\alpha)$ remains strictly positive for all
$|\alpha|$ sufficiently small.
For each $i$ with $\delta\beta_i < 0$, the component
$\beta_i(\alpha)$ reaches zero at
$\alpha_i = \beta_{0,i} / |\delta\beta_i|$.
Define
\[
  \alpha^*
  = \min\bigl\{\,
      \alpha_i : i \in S,\; \delta\beta_i < 0
    \bigr\},
\]
and let $k$ be an index achieving this minimum.
Such an index exists because $\delta\beta$ has at least one
negative component (Step~1).
Note $\alpha^* > 0$ since $\beta_{0,i} > 0$ and
$|\delta\beta_i| > 0$ for each candidate~$i$.

\emph{Step~4: Verify non-negativity.}
Set $\beta^* = \beta(\alpha^*)$.
Then $\beta^*_k = 0$ by construction.
For $j \neq k$:
\begin{itemize}
\item If $\delta\beta_j \geq 0$:
  $\beta^*_j = \beta_{0,j} + \alpha^*\,\delta\beta_j
    \geq \beta_{0,j} > 0$.
\item If $\delta\beta_j < 0$:
  $\alpha^* \leq \alpha_j = \beta_{0,j}/|\delta\beta_j|$
  (since $\alpha^*$ is the minimum), so
  $\beta^*_j
    = \beta_{0,j} - \alpha^*\,|\delta\beta_j|
    \geq \beta_{0,j} - \alpha_j\,|\delta\beta_j| = 0$.
\end{itemize}
Hence $\beta^*_j \geq 0$ for all $j \in S$.

\emph{Step~5: The smaller pair achieves the same value.}
Since $\beta^*_k = 0$, the quadratic form
$\beta^{*\top} H\, \beta^* = \nu_0$ (Step~2)
receives no contribution from index~$k$.
Let $H'$ denote the action matrix for
$(S', \sigma|_{S'})$, where $S' = S \setminus \{k\}$.
The sub-matrix of~$H$ at rows and columns~$S'$
equals~$H'$ (both are defined by the same normals
$n_{\sigma(j)}$, $j \in S'$), so
$(\beta^*|_{S'})^\top H'\, (\beta^*|_{S'}) = \nu_0$.

The restriction $\beta^*|_{S'}$ satisfies:
$\sum_{j \in S'} \beta^*_j\, n_{\sigma(j)} = 0$
(from $N^\top \beta^* = 0$ and $\beta^*_k = 0$),
$\sum_{j \in S'} \beta^*_j\, h_{\sigma(j)} = 1$
(from $\eta^\top \beta^* = 1$ and $\beta^*_k = 0$),
and $\beta^*_j \geq 0$ for $j \in S'$ (Step~4).
Hence $\beta^*|_{S'} \in M(K)$ restricted to~$S'$.
The algorithm evaluates
$\max\, \beta^\top H'\, \beta$ over this constraint set
and therefore finds value $\geq \nu_0$.
\end{proof}

\begin{remark}[Geometric characterisation of the null space]%
\label{rem:nullspace-geometry}
The null space of the system matrix always contains the subspace
$\{(0, \delta\lambda, 0) : N\,\delta\lambda = 0\}$, which has
dimension $4 - \operatorname{rank}(N)$.
In particular, a null-space element with $\delta\beta = 0$
exists if and only if the visited normals
$\{n_{\sigma(i)} : i \in S\}$ do not span~$\mathbb{R}^4$.
(Proof: if $\delta\beta = 0$, the null-space equations reduce to
$N\,\delta\lambda + \eta\,\delta\nu = 0$; multiplying by any
$\beta_0 \geq 0$ with $\eta^\top\!\beta_0 = 1$ gives $\delta\nu = 0$,
leaving $N\,\delta\lambda = 0$.)
The algorithm does not need this characterisation: it suffices
to inspect the null-space vectors from the SVD and check whether
$\delta\beta \neq 0$.
\end{remark}

% Status: NOT YET Jörn-approved.
\begin{proposition}[Approximate dismissal of near-singular systems]%
\label{prop:approximate-dismissal}
Let $(S, \sigma)$ be a pair for which the linear
system~\eqref{eq:linear-system} has a solution
$(\beta_0, \lambda_0, \nu_0)$ with $\beta_{0,i} > 0$ for all
$i \in S$.
Let $(\delta\beta, \delta\lambda, \delta\nu)$ be a unit vector
($\|(\delta\beta, \delta\lambda, \delta\nu)\|_2 = 1$) with
\emph{approximate} residual
\[
  A\begin{pmatrix}\delta\beta\\\delta\lambda\\\delta\nu\end{pmatrix}
  = \begin{pmatrix}r_1\\r_2\\r_3\end{pmatrix},
  \qquad
  \varepsilon
  := \bigl\|(r_1, r_2, r_3)\bigr\|_2,
\]
where $A$ is the system matrix
of~\eqref{eq:linear-system},
and suppose $\delta\beta$ has at least one strictly positive and
one strictly negative component.

Define $\alpha^*$, $k$, and $S' = S \setminus \{k\}$ as in
Lemma~\ref{lem:rank-deficiency-dismissal} (walk to boundary).
Let $\sigma_C > 0$ denote the smallest singular value of the
constraint matrix $C' = [N' \mid \eta'] \in \mathbb{R}^{|S'|\times 5}$
for the smaller pair~$(S', \sigma|_{S'})$.
Then:
\begin{enumerate}
\item
  $\beta^* := \beta_0 + \alpha^*\,\delta\beta$ satisfies
  $\beta^*_k = 0$ and $\beta^*_j \geq 0$ for all $j \in S$.
\item
  The action-value perturbation along the walk satisfies
  \[
    \bigl|\beta^{*\top} H\,\beta^* - \nu_0\bigr|
    \leq \varepsilon\,\alpha^*
       \bigl[
         2\bigl(\|\lambda_0\|_2 + |\nu_0|\bigr) + 3\,\alpha^*
       \bigr]
    =: \Delta_{\mathrm{walk}}.
  \]
\item
  The constraint violations of $\beta^*|_{S'}$ are bounded by
  \[
    \bigl\|{N'}^\top \beta^*|_{S'}\bigr\|_2 \leq \alpha^*\varepsilon,
    \qquad
    \bigl|{\eta'}^\top \beta^*|_{S'} - 1\bigr| \leq \alpha^*\varepsilon.
  \]
\item
  If\/ $\varepsilon$ is small enough that the minimum-norm
  feasibility correction $\Delta\beta$
  satisfies $\|\Delta\beta\|_\infty < \min_{j \in S'} \beta^*_j$
  (so that the corrected point remains non-negative),
  then there exists $\beta' \in M(K)|_{S'}$ with
  \[
    {\beta'}^\top H'\,\beta'
    \geq \nu_0 - \Delta_{\mathrm{walk}} - \Delta_{\mathrm{proj}},
  \]
  where
  $\Delta_{\mathrm{proj}}
     = 2\bigl(\|\beta_0\|_2 + \alpha^*\bigr)
       \,\|H\|_2\,\alpha^*\varepsilon\,/\,\sigma_C$.
\end{enumerate}
Consequently, the maximum of $\beta^\top H'\,\beta$ over
$M(K)|_{S'}$ falls short of $\nu_0$ by at most
$\Delta_{\mathrm{walk}} + \Delta_{\mathrm{proj}}
 = O(\varepsilon)$.
\end{proposition}

\begin{proof}
Conclusion~(1) is identical to
Steps~1--4 of Lemma~\ref{lem:rank-deficiency-dismissal}:
the walk-to-boundary construction and non-negativity argument
are purely algebraic and do not use the exactness of the
null-space relation.

\emph{Step~5${}'$: Action-value perturbation (conclusion~2).}
Expanding the quadratic form at $\beta^*
= \beta_0 + \alpha^*\,\delta\beta$:
\[
  \beta^{*\top} H\,\beta^*
  = \nu_0
    + 2\alpha^*\,\delta\beta^\top H\,\beta_0
    + (\alpha^*)^2\,\delta\beta^\top H\,\delta\beta.
\]
We bound the two inner products using the residual
$(r_1, r_2, r_3)$ in place of the exact null-space
relations of Lemma~\ref{lem:rank-deficiency-dismissal}.

For the cross term: the first-order condition
$H\,\beta_0 = N\,\lambda_0 + \eta\,\nu_0$ gives
\[
  \delta\beta^\top H\,\beta_0
  = (N^\top\delta\beta)^\top \lambda_0
    + (\eta^\top\delta\beta)\,\nu_0
  = r_1^\top \lambda_0 + r_3\,\nu_0.
\]
Hence
$|\delta\beta^\top H\,\beta_0|
  \leq \|r_1\|\,\|\lambda_0\| + |r_3|\,|\nu_0|
  \leq \varepsilon\bigl(\|\lambda_0\| + |\nu_0|\bigr)$.

For the quadratic term: the approximate second block row
$H\,\delta\beta = N\,\delta\lambda + \eta\,\delta\nu + r_2$ gives
\[
  \delta\beta^\top H\,\delta\beta
  = r_1^\top \delta\lambda + r_3\,\delta\nu
    + \delta\beta^\top r_2.
\]
Since $\|\delta\beta\|, \|\delta\lambda\|, |\delta\nu| \leq 1$
(unit vector) and $\|r_1\|, \|r_2\|, |r_3| \leq \varepsilon$:
$|\delta\beta^\top H\,\delta\beta| \leq 3\varepsilon$.

Combining:
$|\beta^{*\top}H\,\beta^* - \nu_0|
  \leq 2\alpha^*\varepsilon(\|\lambda_0\|+|\nu_0|)
     + 3(\alpha^*)^2\varepsilon
  = \Delta_{\mathrm{walk}}$.

\emph{Step~6${}'$: Constraint violation (conclusion~3).}
Since $N^\top\beta_0 = 0$ and $\eta^\top\beta_0 = 1$:
\begin{align*}
  {N'}^\top \beta^*|_{S'}
  &= N^\top\beta^*
   = \alpha^*\,N^\top\delta\beta
   = \alpha^*\,r_1, \\
  {\eta'}^\top \beta^*|_{S'}
  &= \eta^\top\beta^*
   = 1 + \alpha^*\,\eta^\top\delta\beta
   = 1 + \alpha^*\,r_3,
\end{align*}
using $\beta^*_k = 0$ in the restriction.
Hence $\|{N'}^\top\beta^*|_{S'}\| \leq \alpha^*\varepsilon$
and $|{\eta'}^\top\beta^*|_{S'} - 1| \leq \alpha^*\varepsilon$.

\emph{Step~7${}'$: Feasibility correction (conclusion~4).}
The minimum-norm correction
$\Delta\beta$ satisfying
${C'}^\top\Delta\beta = (\alpha^*r_1,\;\alpha^*r_3)^\top$
has
\[
  \|\Delta\beta\|_2
  \leq \frac{\|(\alpha^*r_1,\;\alpha^*r_3)\|_2}{\sigma_C}
  \leq \frac{\alpha^*\varepsilon}{\sigma_C}.
\]
Set $\beta' = \beta^*|_{S'} - \Delta\beta$.
By hypothesis, $\|\Delta\beta\|_\infty < \min_{j \in S'}\beta^*_j$,
so $\beta'_j > 0$ for all $j \in S'$.
By construction, ${N'}^\top\beta' = 0$ and ${\eta'}^\top\beta' = 1$,
so $\beta' \in M(K)|_{S'}$.

The action-value change from the correction is
\begin{align*}
  \bigl|{\beta'}^\top H'\,\beta'
    - (\beta^*|_{S'})^\top H'\,(\beta^*|_{S'})\bigr|
  &\leq 2\,\|\beta^*|_{S'}\|_2\,\|H'\|_2\,\|\Delta\beta\|_2
       + \|H'\|_2\,\|\Delta\beta\|_2^2 \\
  &\leq 2\bigl(\|\beta_0\|_2 + \alpha^*\bigr)
        \|H\|_2\,
        \frac{\alpha^*\varepsilon}{\sigma_C}
       + O(\varepsilon^2)
  = \Delta_{\mathrm{proj}} + O(\varepsilon^2),
\end{align*}
using $\|\beta^*|_{S'}\| \leq \|\beta_0\| + \alpha^*$
and $\|H'\| \leq \|H\|$.
Since the $(S', \sigma|_{S'})$ algorithm maximises over
all of $M(K)|_{S'}$, it finds
value $\geq {\beta'}^\top H'\beta'
       \geq \nu_0 - \Delta_{\mathrm{walk}} - \Delta_{\mathrm{proj}}$.
\end{proof}

\begin{remark}[Algorithmic application]%
\label{rem:approximate-dismissal-algorithm}
In the SVD-based solver, the near-singular direction
$(\delta\beta, \delta\lambda, \delta\nu)$ is the right singular
vector for the smallest singular value~$\sigma_{\min}$,
and $\varepsilon = \sigma_{\min}$.
The condition ``$\delta\beta$ has mixed signs'' holds
whenever $\|\delta\beta\|_2 > \varepsilon / h_{\min}$,
since otherwise
$\sum \eta_i\,\delta\beta_i \geq h_{\min}\|\delta\beta\|_1$
would contradict $|\eta^\top\delta\beta| \leq \varepsilon$.
If $\|\delta\beta\|_2 \leq \varepsilon / h_{\min}$,
the near-singular direction is effectively confined to
the multiplier block and $\beta$ is approximately unique:
the algorithm handles this case without special treatment.

The bound $\Delta_{\mathrm{walk}} + \Delta_{\mathrm{proj}}$
is proportional to $\varepsilon = \sigma_{\min}$ and can be
evaluated numerically for each dropped pair.
\end{remark}
